\makerubrichead{Research Interests}
%
My research develops autonomous \textbf{robotic systems} that can reason, plan, and act in uncertain real-world environments, under physical constraints and safety requirements. While AI has enabled robots to learn complex tasks from data, purely data-driven approaches often lack interpretability, require large datasets, and fail to guarantee safe, reliable behavior. I aim to bridge this gap by combining \textbf{physics-informed} knowledge and structured \textbf{optimization} principles with \textbf{machine learning}, producing planning and control methods that are \textbf{data-efficient}, adaptive, and safety-aware. I couple methodological advances with \textbf{real-world} validation on robotic platforms, targeting practical impact on the next-generation autonomous systems.
%
% The goal of my research is to develop autonomous robotic systems that can reason, plan, and act in uncertain real-world environments, while respecting complex physical constraints, safety requirements, and computational limitations. Recent advances in artificial intelligence have shown how robots can learn from data to perform complex tasks. However, purely data-driven approaches require massive datasets to generalize, lack interpretability, and often fail to provide safe and reliable performance, which is critical for real-world robotics. My research aims to bridge these gaps by exploiting prior physical knowledge and structured optimization principles to develop physics-informed, learning-based robot planning and control methods that are data-efficient, adaptive in changing environments, and compatible with real-time constraints. I combine methodological advances with real-world validation on robotic platforms, aiming for practical impact on the future autonomous systems.
%
% My research lies at the intersection of \textbf{robotics}, \textbf{machine learning}, and \textbf{control} engineering. I develop algorithms for real-time trajectory planning, control, and state estimation, with a focus on mobile ground robots in uncertain, dynamic environments. I like to integrate \textbf{prior knowledge} of the robot dynamics into learning-based methods, such as neural networks, to enhance \textbf{generalization} to unseen scenarios with minimal training data.

\vspace{0.4cm}

% \begin{itemize}
%     \item Trajectory planning and control with autonomous vehicles \vspace{-0.2cm}
%     \item Real-time optimization and model predictive control \vspace{-0.2cm}
%     \item Model-structured neural networks for modeling, control and estimation of dynamical systems \vspace{-0.2cm}
%     \item Learning to control and system identification
% \end{itemize}
